线性模型常被当成入门用的玩具，学完就该换掉。可它至今仍在信贷审批、临床预后、A/B 实验读数和风控规则里当主力，而这些场合对表达力的要求并不低。留下它的理由不在拟合能力，在于每个系数都能指出来源、每份不确定都能量化、每种失效方式都已经被前人清点过并且可以在部署前查出来。可解释和可诊断，是它一直没被换掉的两条腿。

这一单元不按算法名逐个介绍，而是沿着模型遇到的困难往前推。开头是一条几何事实：最小二乘就是把响应向列空间做正交投影，此处还没有任何概率。往上每加一层假设，就多买到一样东西，也多一处可以被数据反驳的地方。

\begin{center}
\begin{tikzpicture}[>=stealth]
  \draw (-1.4,-0.32) rectangle (1.4,0.32);
  \node at (0,0) {\footnotesize 正交投影};
  \draw (-1.4,-1.47) rectangle (1.4,-0.83);
  \node at (0,-1.15) {\footnotesize 三层噪声假设};
  \draw (-1.4,-2.62) rectangle (1.4,-1.98);
  \node at (0,-2.30) {\footnotesize 惩罚项};
  \draw (-1.4,-3.77) rectangle (1.4,-3.13);
  \node at (0,-3.45) {\footnotesize 先验与后验};
  \draw (-1.4,-4.92) rectangle (1.4,-4.28);
  \node at (0,-4.60) {\footnotesize 响应分布与链接};
  \draw[->] (0,-0.32) -- (0,-0.83);
  \draw[->] (0,-1.47) -- (0,-1.98);
  \draw[->] (0,-2.62) -- (0,-3.13);
  \draw[->] (0,-3.77) -- (0,-4.28);
  \node[right,black!60] at (1.7,-0.575) {\footnotesize 换来无偏、方差公式与精确分布};
  \node[right,black!60] at (1.7,-1.725) {\footnotesize 换来稳定，付出偏差};
  \node[right,black!60] at (1.7,-2.875) {\footnotesize 换来完整后验，付出先验的选择};
  \node[right,black!60] at (1.7,-4.025) {\footnotesize 换来新的响应类型，写死均值与方差的关系};
\end{tikzpicture}
\end{center}

\emph{同一条线性骨架，每往下一层就多一份承诺，也多一处可以被数据反驳的地方。}

五篇的分工如下。

\begin{enumerate}
\tightlist
\item
  \hyperref[sec:13-Machine-Learning/02-Linear-Models/01]{``线性回归从投影到概率模型''}：从两套同小区同户型却成交价相差一千二的房子出发，说清最小二乘的几何唯一性、正规方程为什么不该拿来求解，以及拟合成立、统计推断成立与因果解释成立各自需要什么条件。
\item
  \hyperref[sec:13-Machine-Learning/02-Linear-Models/02]{``正则化从 Ridge 到 Lasso''}：特征近似共线时数据无法确定系数。Ridge 同时结统计与数值两笔账，Lasso 的稀疏来自范数在原点的不可导，而共线时它选中谁带有任意性。
\item
  \hyperref[sec:13-Machine-Learning/02-Linear-Models/03]{``从二分类到 Softmax 分类''}：响应变成类别后，线性结构该安在对数几率上。系数是几率的倍数而不是概率的差值；数据可分时最大似然跑向无穷，得靠惩罚拉回有限。
\item
  \hyperref[sec:13-Machine-Learning/02-Linear-Models/04]{``Bayesian 线性回归把系数不确定性传到预测''}：先验精度与数据精度按方向相加。Ridge 只用了后验的众数，而预测方差分成观测噪声与参数不确定两段，只有后者会随样本缩小。
\item
  \hyperref[sec:13-Machine-Learning/02-Linear-Models/05]{``GLM 用链接函数连接均值与线性预测子''}：把响应分布、链接函数和线性预测子拆成三个可以分别指定的部件。链接函数的选择是建模决定，它连带改变系数含义、外推形状与不确定性走向。
\end{enumerate}

\subsection{怎么读这一章}

只走主干的话，前三篇按顺序读：它们分别给出投影与噪声假设、稳定性与复杂度、以及从分数到概率的连接。后两篇是两条扩展，一条把点估计换成完整后验，一条把响应分布换掉。第一遍不必记住任何闭式解，能从假设重建目标函数、能说出系数稳定与预测有效与因果解释为什么不是同一件事，这一层就够了。

有几个问题值得在每一篇里反复问：模型把什么当固定、什么当随机？目标函数是从生成假设推出来的还是随手挑的？参数能否被数据唯一识别？数值解可靠吗？验证集上的好成绩能迁移到真正关心的人群吗？这些问题在深层模型里同样存在，只是那时更难看清答案。

读完可以接\hyperref[ch:112]{``降维与低秩建模：数据为何看似高维''}，去问特征表示本身能不能用更少的方向解释，而不是急着换一个更复杂的预测器。第四篇里的 Gaussian 共轭更新会在\hyperref[sec:13-Machine-Learning/06-Probabilistic-Models/07]{``共轭先验让 Bayesian 更新保持同一形式''}与\hyperref[sec:13-Machine-Learning/06-Probabilistic-Models/06]{``指数族把归一化、矩与似然连接起来''}里推广；把系数的解释推到干预层面，则要进\hyperref[ch:120]{``因果推断：预测世界不等于改变世界''}。

线性模型的表达力早就被超过了，把它读薄不难，读厚才有价值。这一单元想留下的判断是：一个模型值不值得信，取决于它的假设有没有被写在明面上、有没有对应的检查手段，而不取决于它有多复杂。
